% ---
% filename : nibt_installations_20260517_180_protection_neutre_reduit.tex
% id : nibt_installations_20260517_180
% title : "Protection du conducteur neutre de section réduite"
% domain : "nibt"
% subdomain : "Neutre"
% tags : ["nibt", "neutre", "section", "surcharge", "dimensionnement", "dt"]
% difficulty : 2
% time_solve : 15
% has_octave : true
% octave_files : ["nibt_installations_20260517_180_protection_neutre_reduit.m"]
% author : "om"
% ---
\documentclass[a4paper,11pt,answers,addpoints]{exam}
% ---------------------------------------------------------
% LANGUE, ENCODAGE, FONTS
% ---------------------------------------------------------
\usepackage[utf8]{inputenc}
% ---------------------------------------------------------
% GÉOMÉTRIE ET MISE EN PAGE
% ---------------------------------------------------------
\usepackage[left=1.7cm,right=1.5cm,top=2cm,bottom=1cm]{geometry}
\usepackage{microtype}
\usepackage{float}
\usepackage{needspace}

% ---------------------------------------------------------
% MATHÉMATIQUES
% ---------------------------------------------------------
\usepackage{amsmath, amssymb, bm, esvect}

\newcommand{\V}{\overrightarrow}
\def\vect#1{\overrightarrow{\strut#1}}
\providecommand{\Degrees}[0]{\ensuremath{^\circ}}

% ---------------------------------------------------------
% UNITÉS ET GRANDEURS (SIunitx)
% ---------------------------------------------------------
\usepackage{siunitx}
\sisetup{output-decimal-marker={,}}
\DeclareSIUnit \var {var}
\DeclareSIUnit \VA {VA}
\DeclareSIUnit \kVA {kVA}
\DeclareSIUnit[quantity-product={}] \degree{\text{\textdegree}}

% Permet la césure dans les nombres avec unités (évite les Underfull \hbox)
%\AllowBreakAtQQ

% ---------------------------------------------------------
% GRAPHIQUES : TikZ + CircuitikZ (sans tkz-fct qui cause des erreurs spath3)
% ---------------------------------------------------------
\usepackage{tikz}
\usetikzlibrary{
	arrows.meta,
	intersections,
	decorations.pathreplacing,
	positioning
}

\usepackage[european,straightvoltages,EFvoltages]{circuitikz}
\usepackage{pgfplots}
\pgfplotsset{compat=1.12}

% Symbole étoile-triangle personnalisé
\newcommand{\wye}{%
	\mathbin{\tikz[x=1ex,y=1ex]{%
			\draw[line width=.1ex] (0,0)--(45:1)--++(-45:1) (45:1)--++(0,1);
	}}%
}

% ---------------------------------------------------------
% TABLEAUX
% ---------------------------------------------------------
\usepackage{tabularx}
\usepackage{tabu}

% ---------------------------------------------------------
% HYPERLIENS
% ---------------------------------------------------------
\usepackage[colorlinks=true,
menucolor=red,
breaklinks=true,
urlcolor=blue,
linkcolor=blue]{hyperref}

% ---------------------------------------------------------
% COMMANDES PERSONNELLES
% ---------------------------------------------------------
\newcommand\nomauteur{bdminasmorgul@protonmail.com}
\newcommand\dateTe{18.05.2026}
\newcommand\brancheTe{DT}

% ---------------------------------------------------------
% STYLE DE L'EXERCICE
% ---------------------------------------------------------
\pagestyle{headandfoot}
\firstpageheadrule
\runningheadrule

% En-tête avec largeur fixe pour éviter les dépassements
\firstpageheader{\brancheTe\ (\nomauteur)}{Élève : }{\makebox[3.5cm][r]{\dateTe\ \thepage/\numpages}}
\runningheader{\brancheTe\ (\nomauteur)}{Élève : }{\makebox[3.5cm][r]{\dateTe\ \thepage/\numpages}}

% Pied de page
\newcommand{\totalpointtts}{%
	\begin{tikzpicture}[remember picture, overlay]
		\node[anchor=south west, inner sep=0pt, outer sep=0pt]
		at ([xshift=0.5cm,yshift=0.5cm]current page.south west)
		{\rotatebox{90}{Total des points : \numpoints}};
	\end{tikzpicture}%
}

\firstpagefooter{\hspace*{\fill}\totalpointtts}{}{}
\runningfooter{\hspace*{\fill}\totalpointtts}{}{}

% Réglages exam
\printanswers
\addpoints
\pointsinmargin
\bracketedpoints

% ---------------------------------------------------------
% LISTINGS (code)
% ---------------------------------------------------------
\usepackage{xcolor}
\usepackage{listings}

\lstdefinestyle{octavestyle}{
	language=Matlab,
	basicstyle=\ttfamily\small,
	keywordstyle=\color{blue},
	commentstyle=\color{green!50!black},
	stringstyle=\color{red},
	stepnumber=1,
	frame=single,
	breaklines=true,
	breakatwhitespace=true,
	tabsize=4
}

\usepackage{enumitem}
\setlist[enumerate,1]{label=\alph*), ref=\alph*)}
\newenvironment{explication}{%
	\par\noindent\textbf{Explication. }\itshape
}{\par\normalfont}

\begin{document}
	\unframedsolutions
	\pointsinmargin
	\bracketedpoints
	
\section*{Préparation DT PQ CFC ELMO,IELE,PELE}
\begin{questions}
\question
Comment réaliser le dimensionnement des colonnes montantes pour un immeuble locatif. Dans certains circuits triphasés, vous envisagez d'utiliser un conducteur neutre de section inférieure à celle des conducteurs de phase.

\begin{parts}
	\part[1] Quelle disposition doit être prise lorsque la section du conducteur neutre ($S_N$) est inférieure à celle des conducteurs de phase ($S_L$)~?
	\part[3] Selon la NIBT 5.2.4.3, quelle est la section minimale absolue ($S_{N,min}$) pour un conducteur neutre réduit en cuivre, et quel pourcentage minimal de la section de phase doit-il respecter~ ?
	\part[3] Dans quelle condition spécifique liée aux courants harmoniques est-il strictement interdit de réduire la section du conducteur neutre (NIBT 5.2.3.4.2)~ ?
	\part[3] Si un conducteur neutre réduit est utilisé, doit-on obligatoirement prévoir un dispositif de détection de surintensité spécifique pour le neutre s'il est déjà protégé contre les courts-circuits par le dispositif des phases et que le courant de service est faible~ ?
\end{parts}

\begin{solution}
	\begin{itemize}
		\item \textbf{Disposition impérative} : Il faut impérativement s'assurer que le conducteur de neutre ne puisse pas être surchargé (NIBT 4.3.1.2.2).
		\item \textbf{Limites de réduction} : La section réduite du neutre doit être au moins égale à \qty{50}{[\percent]} de celle de la phase, avec un minimum de \qty{16}{[\milli\metre\squared]} pour le cuivre (NIBT {5.2.4.3}).
		\item \textbf{Interdiction (Harmoniques)} : La réduction est interdite si le taux d'harmoniques (notamment de rang 3) est supérieur à \qty{15}{[\percent]}, car le courant de neutre peut alors dépasser celui des phases (NIBT {5.2.3.4.2}).
		\item \textbf{Protection contre les surintensités} : Il n'est pas nécessaire de prévoir une détection spécifique si le neutre est protégé contre les courts-circuits par le dispositif des phases et que le courant maximal prévisible en service normal est inférieur au courant admissible du neutre réduit (NIBT {4.3.1.2.2.1}).
	\end{itemize}
	

	\begin{align*}
		S_L > \qty{16}{[\milli\metre\squared]} &\implies S_N \geq \dfrac{S_L}{2} \text{ et } S_N \geq \qty{16}{[\milli\metre\squared]} \\[6pt]
		I_{harmoniques} > \qty{15}{[\percent]} &\implies S_N = S_L \\[6pt]
		\text{Condition de sécurité} &: I_{N,max} < I_{z,N}
	\end{align*}
	
	
\end{solution}
\end{questions}
\end{document}
